\section{Regional energy in a close candidate space}

The regional excitation energy of a global ground vector is controlled
by the bonds crossing the region boundary. An orthogonal projection
preserves a constant energy bound when its loss is small enough. These
estimates support the energy-estimation step of~\cite{Arad2016RigorousRG}.

\begin{lemma}[Regional ground energy and its projected image]
    \label{lem:regional_ground_energy_bound}
    \lean{TNLean.MPS.AreaLaw.exists_regional_ground_energy_bound}
    \leanok
    Let $H=\sum_i h_i$ be a chain Hamiltonian with $0\le h_i\le I$
    supported on sites $i,i+1$, and let the local dimension be positive.
    Fix an interval $J$. Let $A_J$ be the sum of the bonds wholly inside
    $J$, and let $G$ be the full ground space of $H$. The regional minimum
    $e_J$ is attained, and the operator $C=A_J-e_JI$ satisfies
    \[
       C\ge0,\qquad \|C\|\le L,\qquad
       \operatorname{Re}\langle Cg,g\rangle\le2\|g\|^2
       \quad(g\in G).
    \]
    These data are chosen before every candidate subspace $U$ and every
    $0\le\delta\le[100(L+1)]^{-1}$. If
    $\|P_Ug\|^2\ge(1-\delta)\|g\|^2$ for every $g\in G$, then
    \[
       \operatorname{Re}\langle Cv,v\rangle\le3\|v\|^2
       \quad(v\in P_U G).
    \]
    There is also a nonzero $v\in U$ satisfying this energy bound.
    The statement includes empty intervals and chain endpoints, and
    makes no assumption of a gap or a unique ground vector.
\end{lemma}
\begin{proof}\leanok
    \uses{lem:alvv_single_region_residuals}
    Let $e_0,e_1,e_2$ be the actual minima on the three disjoint regions,
    with $e_1=e_J$. The original regional ground data give
    \[
       0\le\lambda_0-(e_0+e_1+e_2)\le2,
       \qquad C\le H-(e_0+e_1+e_2)I,
    \]
    where $\lambda_0$ is the least eigenvalue of $H$. The upper bound
    two is the norm bound for the two crossing bonds. Apply this order
    inequality on $G$ to obtain the first energy bound. Positivity of
    $e_J$ and the bond count give $0\le C\le LI$ and hence $\|C\|\le L$.

    For $v=P_Ug$, positivity at $g-8(v-g)$ gives
    \[
       \operatorname{Re}\langle Cv,v\rangle
       \le\frac98\operatorname{Re}\langle Cg,g\rangle
          +9\operatorname{Re}\langle C(v-g),v-g\rangle.
    \]
    Pythagoras gives $\|v-g\|^2\le\delta\|g\|^2$, and the assumed
    projection inequality bounds the projected norm from below. Thus
    the energy on the image is at most
    $(9/4+9L\delta)/(1-\delta)$ per unit squared norm, which is at
    most three. Project an actual unit ground eigenvector to obtain
    the nonzero candidate vector; $\delta<1$ ensures it does not vanish.
\end{proof}

The compressed regional minimum and the recursive update of its energy
estimate are subsequent conclusions.

\begin{lemma}[Energy estimates under joining disjoint regions]
    \label{lem:disjoint_regional_energy_merge}
    \lean{TNLean.MPS.AreaLaw.exists_disjoint_regional_energy_merge}
    \leanok
    Let $A,B\ge0$ be supported on disjoint site sets, and let
    $0\le Y\le I$. There are attained minimum energies $e_A,e_B,e_P$
    of $A,B,A+B+Y$, respectively, with a common unit minimizing vector
    for $A$ and $B$. They satisfy
    \[
       e_A+e_B\le e_P\le e_A+e_B+1.
    \]
    Consequently, if $p_A\in[e_A-3,e_A]$ and
    $p_B\in[e_B-3,e_B]$, then
    \[
       p_A+p_B\in[e_P-7,e_P].
    \]
    No gap, unique ground vector, or minimum-energy witness is assumed.
\end{lemma}
\begin{proof}\leanok
    \uses{thm:disjoint_regional_clipping}
    Choose a common regional minimizing vector using disjoint supports.
    The centered operator $X=A-e_AI+B-e_BI$ is positive and vanishes
    on that vector. The minimum $E$ of $X+Y$ therefore belongs to $[0,1]$.
    Adding back the two regional energies gives
    $e_P=e_A+e_B+E$. The estimate for $p_A+p_B$ follows by addition.
    This is auxiliary to the energy-estimation argument in
    \cite[Figure \texttt{fig:algo1}, lines 1627--1632]{Arad2016RigorousRG}.
\end{proof}

\begin{lemma}[Actual adjacent-interval energy estimates]
    \label{lem:adjacent_interval_energy_merge}
    \lean{TNLean.MPS.AreaLaw.exists_adjacent_interval_energy_merge}
    \leanok
    Let $H=\sum_i h_i$ be a nearest-neighbour chain with $0\le h_i\le I$
    and positive local dimension. Fix $0\le l\le m\le r\le L$.
    Write $A,B,P$ for the internal bond sums on $[l,m)$, $[m,r)$,
    and $[l,r)$, respectively. Their attained minima satisfy
    \[
       e_A+e_B\le e_P\le e_A+e_B+1.
    \]
    A common unit vector attains the two child minima. If
    $p_A\in[e_A-3,e_A]$ and $p_B\in[e_B-3,e_B]$, then
    \[
       p_A+p_B\in[e_P-7,e_P].
    \]
    Empty children, empty parents and chain endpoints are included.
    No gap, energy witness, or separate regional decomposition is assumed.
\end{lemma}
\begin{proof}\leanok
    \uses{lem:disjoint_regional_energy_merge}
    Partition the parent bonds into the two internal child families and
    the bonds crossing the cut at $m$ while remaining internal to the parent.
    There is at most one such bond. Its sum $Y$ satisfies $0\le Y\le I$
    and $P=A+B+Y$. The two child operators are supported on disjoint site
    sets, so the preceding lemma applies.
\end{proof}
